\documentclass[11pt,a4paper]{article}
\usepackage[utf8]{inputenc}
\usepackage[T1]{fontenc}
\usepackage[spanish]{babel}
\usepackage{amsmath, amssymb}
\usepackage{enumitem}
\usepackage{geometry}
\geometry{margin=2.5cm}

\begin{document}

\subsection*{1. ¿Se hizo una simulación Monte Carlo?}

\textbf{Sí, en su variante de Muestreo por Aceptación y Rechazo (\textit{Acceptance-Rejection Sampling}).}

En lugar de calcular analíticamente las probabilidades conjuntas sobre las $32,468,436$ combinaciones con múltiples restricciones no lineales simultáneas (pesos empíricos, rango de suma, paridad y no consecutividad), el código implementa un \textbf{experimento estocástico repetido}:

\begin{enumerate}
    \item \textbf{Distribución instrumental (Propuesta):} \\
    Se define un vector de probabilidades $\mathbf{p} = (p_1, p_2, \dots, p_{56})$ normalizado ($\sum p_i = 1$) que distorsiona la distribución uniforme con base en la evidencia histórica (pesos por frecuencia o rezago).
    
    \item \textbf{Generación de variables pseudoaleatorias:} \\
    Se ejecuta una selección sin reemplazo:
    $$X^{(k)} \sim \text{Multinomial}(\mathbf{p}, 6)$$
    
    \item \textbf{Criterio de Aceptación/Rechazo:} \\
    Se evalúa si la muestra $X^{(k)}$ satisface simultáneamente las funciones indicadoras de la región objetivo $\Omega_{\text{objetivo}}$:
    $$\mathbb{I}_{\text{suma}}\big(X^{(k)}\big) \times \mathbb{I}_{\text{paridad}}\big(X^{(k)}\big) \times \mathbb{I}_{\text{adyacencia}}\big(X^{(k)}\big) = 1$$
    \begin{itemize}
        \item Si el producto es $1$, el estado se acepta como candidato viable.
        \item Si es $0$, el estado se descarta (\textit{rechazo}) y el contador de iteraciones $k \to k+1$ continúa hasta converger al número solicitado de combinaciones.
    \end{itemize}
\end{enumerate}

\textit{(Nota teórica: Es un método Monte Carlo de muestreo condicional, no una simulación masiva de horizonte temporal largo para calcular probabilidades de ruina o valor esperado acumulado).}

\medskip
\hrule
\medskip

\subsection*{2. Índice Metodológico para el Análisis de Espacios Muestrales Futuros}

Este índice sirve como marco de referencia formal para modelar cualquier juego de azar o sistema combinatorio finito discreto (Melate, Revancha, Chispazo, Tris, etc.):

\begin{enumerate}[label=\textbf{\arabic*.}]
    \item \textbf{TOPOLOGÍA DEL ESPACIO MUESTRAL}
    \begin{enumerate}[label=\textbf{1.\arabic*}]
        \item Definición del alfabeto o conjunto base $S = \{1, 2, \dots, N\}$.
        \item Dimensión de la tupla ($k$ lugares sin reemplazo / con reemplazo).
        \item Cardinalidad teórica del espacio ($|\Omega| = C(N, k)$ ó $P(N, k)$).
        \item Medida de probabilidad base (Espacio equiprobable de Laplace: $P(\omega) = \frac{1}{|\Omega|}$).
    \end{enumerate}

    \item \textbf{PROYECCIÓN Y PARTICIÓN MEDIANTE VARIABLES ALEATORIAS}
    \begin{enumerate}[label=\textbf{2.\arabic*}]
        \item Variable de Suma Total: $S(\omega) = \sum r_i$ y su función de masa inducida.
        \item Variable de Paridad: Número de éxitos en distribución hipergeométrica/binomial.
        \item Variable de Dispersión / Amplitud: $R(\omega) = \max(\omega) - \min(\omega)$.
        \item Detección de subconjuntos de alta densidad (Regiones hipervolumétricas con mayor $\frac{|A|}{|\Omega|}$).
    \end{enumerate}

    \item \textbf{CONTRASTE EMPÍRICO CONTRA EL HISTORIAL $D$}
    \begin{enumerate}[label=\textbf{3.\arabic*}]
        \item Vector de frecuencias marginales observadas $\mathbf{f} = (f_1, \dots, f_N)$.
        \item Prueba de hipótesis de bondad de ajuste (Chi-cuadrada: $H_0: p_i = \frac{1}{N}$).
        \item Función de supervivencia y rezago (Distribución geométrica del tiempo entre éxitos).
        \item Matriz de co-ocurrencia (Frecuencia empírica de pares y ternas observadas).
    \end{enumerate}

    \item \textbf{INFERENCIA Y MODELADO PROBABILÍSTICO}
    \begin{enumerate}[label=\textbf{4.\arabic*}]
        \item Inferencia Bayesiana: Actualización de \textit{priors} mediante verosimilitud histórica.
        \item Cadenas de Markov de primer orden: $P(X_{t+1} = j \mid X_t = i)$.
        \item Ponderación de entropía máxima sujeta a restricciones empíricas.
        \item Condicionamiento del espacio: Restricción a la subálgebra $\Omega^* \subset \Omega$.
    \end{enumerate}

    \item \textbf{MÉTODOS ESTOCÁSTICOS DE EXTRACCIÓN (MONTE CARLO)}
    \begin{enumerate}[label=\textbf{5.\arabic*}]
        \item Muestreo por Aceptación y Rechazo (\textit{Rejection Sampling}).
        \item Muestreo de Importancia (\textit{Importance Sampling} para eventos raros).
        \item Algoritmo Metropolis-Hastings / MCMC (para espacios de estados complejos).
        \item Simulación de Monte Carlo cruda (Convergencia por Ley Fuerte de Grandes Números).
    \end{enumerate}

    \item \textbf{EVALUACIÓN Y VALIDACIÓN DEL MODELO}
    \begin{enumerate}[label=\textbf{6.\arabic*}]
        \item \textit{Backtesting}: Rendimiento del modelo evaluado contra sorteos pasados ($t-n$ hasta $t$).
        \item Entropía de Shannon del generador (Medición del grado de aleatoriedad vs. determinismo).
        \item Tasa de aceptación del algoritmo (Eficiencia computacional del muestreador).
    \end{enumerate}
\end{enumerate}

\end{document}